\chapter{TRABALHOS RELACIONADOS}

Este capítulo apresenta uma revisão da literatura pertinente ao escopo deste trabalho, buscando fundamentar as escolhas metodológicas adotadas na construção do laboratório experimental e fornecer bases teóricas para justificar os resultados empíricos obtidos. Os estudos aqui selecionados discutem a eficácia e o custo computacional de algoritmos de \textit{Machine Learning} e \textit{Deep Learning} aplicados a Sistemas de Detecção de Intrusão (IDS), os severos desafios de generalização de modelos frente a ameaças inéditas ou classes omitidas, e as recentes tendências de otimização utilizando arquiteturas híbridas e filtros de redução de dimensionalidade. 

\section{\textit{Network Intrusion Detection Systems Design: A Machine Learning Approach}}

Um estudo pertinente para a fundamentação deste trabalho é a pesquisa conduzida por Neto e Gomes \cite{neto_gomes}, intitulada \textit{"Network Intrusion Detection Systems Design: A Machine Learning Approach"}. Os autores propõem uma abordagem sistemática para apoiar a tomada de decisão na seleção de algoritmos de \textit{Machine Learning} no projeto de Sistemas de Detecção de Intrusão (\textit{IDS}).

Os autores elegem o \textit{dataset} \textit{CICIDS2017}, justificando que bases tradicionais (como o \textit{KDD99}) encontram-se defasadas frente à diversidade de tráfego e à complexidade das ameaças modernas. O \textit{pipeline} de pré-processamento descrito por eles inclui a remoção de valores infinitos (considerados ruídos) e a aplicação da normalização \textit{Min-Max} para adequar a escala dos dados numéricos antes do treinamento.

No que tange à seleção dos modelos, o estudo correlato avaliou oito algoritmos distintos de aprendizado de máquina. Os resultados demonstraram que o \textit{Random Forest} e o \textit{Multi-layer Perceptron} (\textit{MLP}), uma arquitetura de Rede Neural Profunda (\textit{DNN}), figuraram entre os algoritmos com os melhores desempenhos, atingindo métricas de Precisão, \textit{Recall} e \textit{F1-Score} na faixa de 98\% a 99\%. Adicionalmente, os autores destacam o tempo de classificação em dados não vistos como um fator crítico para a viabilidade de um \textit{IDS} no mundo real.

O trabalho de Neto e Gomes apresenta ainda técnicas adicionais de interesse. Para lidar com a alta dimensionalidade do \textit{CICIDS2017}, os autores aplicaram a técnica de seleção de características \textit{Mean Decrease in Impurity} (\textit{MDI}), utilizando o próprio algoritmo \textit{Random Forest} para reduzir o conjunto de entrada de 78 para apenas 51 \textit{features}. Essa abordagem demonstrou ser eficaz para reduzir o custo de treinamento sem perda significativa de desempenho.

Outro ponto de destaque do artigo é a análise comparativa de técnicas de reamostragem para avaliação dos modelos. Os autores demonstraram que o uso da validação cruzada estratificada (\textit{Stratified 10-Fold Cross-Validation}) melhora consideravelmente a estabilidade dos resultados e lida melhor com o severo desbalanceamento de classes inerente ao \textit{CICIDS2017}. Essa técnica garante que a mesma proporção de cada classe de ataque seja rigorosamente mantida durante as divisões de treinamento e teste. Essa evidência demonstra que a estratificação é especialmente relevante para classes com baixo suporte em conjuntos severamente desbalanceados.

\section{\textit{A Comparative Analysis of Signature-Based and Anomaly-Based Intrusion Detection Systems}}

Para compreender os desafios inerentes à adoção de algoritmos de aprendizado de máquina em ambientes de rede reais, a recente pesquisa de Ravindran, Ojha e Kamboj \cite{ravindran2025}, intitulada \textit{"A Comparative Analysis of Signature-Based and Anomaly-Based Intrusion Detection Systems"}, oferece uma boa base teórica. Os autores conduzem uma análise comparativa abrangente focada em métricas operacionais, avaliando o equilíbrio entre a precisão na detecção, o custo computacional, a taxa de falsos positivos e a adaptabilidade a novas ameaças (ataques \textit{zero-day}).

O estudo correlato enfatiza que sistemas baseados em anomalia possuem uma vantagem inerente na adaptabilidade a ameaças inéditas, mas sofrem criticamente com a alta sobrecarga computacional e elevadas taxas de Falsos Positivos. Esses achados apontam para a importância de avaliar, além das métricas de classificação, o tempo de inferência e a análise detalhada da matriz de confusão em comparações entre modelos de IDS.

A literatura aponta ainda caminhos promissores para a evolução de modelos de IDS. O primeiro deles diz respeito à interpretabilidade dos resultados. Enquanto o \textit{Random Forest} possui maior interpretabilidade natural, as Redes Neurais Profundas atuam predominantemente como modelos "caixa-preta". Os autores destacam que a falta de transparência dificulta a confiança dos analistas de segurança. Como melhoria futura, sugere-se a integração de mecanismos de explicabilidade (\textit{Explainable AI}, XAI), como SHAP (\textit{SHapley Additive exPlanations}) ou \textit{LIME}, para tornar as decisões da DNN interpretáveis aos analistas de segurança.

Por fim, o estudo ressalta uma forte tendência na adoção de arquiteturas de IDS Híbridas para lidar com a complexidade dos ataques modernos. A proposta consiste em empregar a detecção por assinaturas como primeiro filtro para ameaças conhecidas e reservar o processamento dos modelos de \textit{Machine Learning} exclusivamente para o tráfego suspeito ou desconhecido.

\section{Detecção de Ataques Web: Explorando Redes Neurais Recorrentes com Redutor de Dimensionalidade}

A preocupação com o alto custo computacional inerente aos modelos de \textit{Deep Learning} é um tema central na literatura de segurança, sendo profundamente investigado por Rego e Nunes \cite{rego_nunes} no artigo \textit{"Detecção de Ataques Web: Explorando Redes Neurais Recorrentes com Redutor de Dimensionalidade"}. Os autores partem da premissa de que as pesquisas com redes neurais têm focado exaustivamente no aumento do desempenho preditivo, negligenciando a eficiência em termos de tempo de detecção. 
Para solucionar o gargalo do custo computacional das redes neurais, os autores propõem o método \textit{BLOOM-RNN}, uma arquitetura de detecção em dois estágios. No primeiro estágio, os autores empregam um Filtro de Bloom, uma estrutura de dados probabilística e altamente eficiente, alimentado exclusivamente com assinaturas de requisições de rede normais (benignas). O objetivo desse filtro é atuar como um redutor de dimensionalidade em nível de volume de tráfego, identificando e descartando requisições sabidamente seguras em frações de milissegundos. Apenas as requisições que o filtro não reconhece como benignas são encaminhadas para o segundo estágio, onde a Rede Neural realiza a classificação profunda.

Essa arquitetura de dois estágios é especialmente adequada para conjuntos com alto desbalanceamento como o \textit{CICIDS2017}, onde a classe \textit{BENIGN} representa a esmagadora maioria do tráfego. Ao filtrar o tráfego normal precocemente, a carga sobre a DNN seria drasticamente reduzida, e seu poder de extração de características seria direcionado exclusivamente para a identificação de ataques.

Ademais, o trabalho de Rego e Nunes reforça a importância da etapa de extração de características (\textit{Feature Extraction}). Os autores demonstraram que é possível obter métricas de acurácia acima de 95\% fornecendo apenas 9 características cruciais para a rede neural, em contraste com os 78 atributos presentes no \textit{CICIDS2017}. Essa evidência reforça a recomendação de combinar seleção de \textit{features} com filtragem de primeiro estágio para viabilizar a implantação de DNNs em sistemas de detecção em tempo real.

\section{\textit{HSAE: A Hybrid Unsupervised Autoencoder for Zero-Day Attack Detection in Realistic Environments}}

A limitação de generalização de modelos supervisionados frente a ataques desconhecidos é amplamente discutida por Silva et al. \cite{silva_hsae} no artigo \textit{"HSAE: A Hybrid Unsupervised Autoencoder for Zero-Day Attack Detection in Realistic Environments"}. Os autores argumentam que modelos supervisionados dependem de dados rotulados e falham diante de ataques \textit{zero-day}, pois só reconhecem classes representadas no treinamento.

Como proposta de mitigação para essas limitações e direcionamento para trabalhos futuros, o estudo correlato introduz o modelo HSAE (\textit{Hybrid Scoring Autoencoder}), uma arquitetura híbrida, leve e computacionalmente eficiente. O HSAE utiliza um paradigma não supervisionado, sendo treinado exclusivamente com tráfego de rede benigno. A detecção de anomalias ocorre pela combinação do erro de reconstrução dos dados no \textit{Autoencoder} com uma função auxiliar de pontuação (\textit{anomaly score}), permitindo separar o comportamento normal do anômalo com alta precisão e baixo tempo de inferência.

Treinar exclusivamente com tráfego benigno elimina a dependência de rótulos de ataque e permite sinalizar qualquer desvio como potencial ameaça, sem restrições quanto ao tipo de ataque.

Adicionalmente, os métodos não supervisionados costumam sofrer com altas taxas de falsos positivos devido à grande variabilidade do tráfego legítimo. Para contornar esse problema, os autores substituíram o uso de limiares de decisão estáticos por um mecanismo de otimização dinâmica baseado na métrica de \textit{Equal Error Rate} (EER).

\section{\textit{Machine Learning in Network Intrusion Detection: A Cross-Dataset Generalization Study}}

A capacidade de um modelo de Inteligência Artificial manter seu desempenho ao analisar dados não vistos é um requisito fundamental para a implantação de Sistemas de Detecção de Intrusão no mundo real. O estudo de Cantone, Marrocco e Bria \cite{cantone2024}, intitulado \textit{"Machine Learning in Network Intrusion Detection: A Cross-Dataset Generalization Study"}, investiga a capacidade de generalização de modelos de aprendizado de máquina em cenários cruzados. Os autores realizaram uma extensa análise cruzada entre diferentes \textit{datasets} de segurança e constataram que, embora os modelos de \textit{Machine Learning} alcancem desempenho quase perfeito quando treinados e testados no mesmo conjunto de dados, a acurácia despenca para níveis próximos à aleatoriedade ao enfrentar dados de redes ou cenários distintos.

Os autores argumentam que os \textit{datasets} públicos atuais possuem classes de ataques excessivamente homogêneas, o que leva os algoritmos a decorarem correlações específicas e artificiais do \textit{dataset} (sobreajuste) em vez de aprenderem o padrão comportamental verdadeiro do ataque. 

Além do sobreajuste, o estudo revela que o desempenho insatisfatório na generalização está intrinsecamente ligado a falhas estruturais no próprio \textit{dataset} CICIDS2017. A literatura aponta que o conjunto original apresenta problemas críticos de integridade, incluindo duplicação de características, erros matemáticos no cálculo de fluxos, falhas na detecção de protocolos e, principalmente, incorreções severas na rotulagem dos ataques. Para contornar essas falhas, os autores recorreram às bases \textit{LycoS-IDS2017} e \textit{LycoS-Unicas-IDS2018}, versões com características reextraídas e rótulos corrigidos a partir dos pacotes brutos originais, confirmando que a qualidade do \textit{dataset} impacta diretamente a capacidade de generalização dos modelos.

Um ponto metodológico relevante é que os autores não utilizaram todas as características disponíveis de forma indiscriminada. No pré-processamento do CICIDS2017, os autores removeram campos como \textit{Flow ID}, \textit{Source IP}, \textit{Source Port}, \textit{Destination IP} e \textit{Protocol}, por representarem identificadores diretos ou quase diretos dos fluxos de rede. Também descartaram o \textit{Timestamp}, que poderia induzir o classificador a aprender relações temporais específicas da coleta, e uma das colunas \textit{Fwd Header Length}, devido à redundância presente no conjunto original. Essa remoção é importante porque tais atributos podem permitir que o modelo memorize particularidades artificiais do ambiente de captura, elevando a acurácia interna, mas reduzindo sua capacidade de generalizar para outras redes.

Os autores também demonstraram que arquiteturas de aprendizado tendem a sofrer com a redundância de amostras e com a alta dimensionalidade do \textit{CICIDS2017} original. Em seus experimentos, a aplicação de seleção de características por \textit{Minimum Redundancy Maximum Relevance} (mRMR) e o uso de subconjuntos reduzidos, de três a cinco características, forçaram os modelos a aprenderem padrões mais gerais e diminuíram o \textit{overfitting} cruzado.

Os achados dos autores sugerem que abordagens como o \textit{Random Forest} tendem a apresentar comportamento mais estável frente a limitações de dados, e que parte da dificuldade em classificar certas ameaças em \textit{datasets} públicos reflete diretamente a qualidade dos dados subjacentes.

\section{Cenários do Experimento}

A análise da literatura permite identificar três linhas de investigação com potencial de ampliar os resultados de experimentos de detecção de intrusão baseados em aprendizado de máquina.

A primeira linha aborda a qualidade dos dados. Cantone et al. \cite{cantone2024} demonstraram que o \textit{CICIDS2017} original apresenta erros de rotulagem, duplicação de características e inconsistências no cálculo de fluxos, e que versões corrigidas como o \textit{LycoS-IDS2017} melhoram a capacidade de generalização dos classificadores.

A segunda linha aborda a alta dimensionalidade do conjunto de características. Tanto Neto e Gomes \cite{neto_gomes} quanto Cantone et al. \cite{cantone2024} mostraram que a seleção de características reduz redundância e pode melhorar o desempenho de classificadores neurais. A técnica \textit{Mean Decrease in Impurity}, combinada com pesos de classe balanceados, é apontada como estratégia promissora para conjuntos com alto desbalanceamento de classes.

A terceira linha trata da limitação de generalização frente a ataques desconhecidos. Silva et al. \cite{silva_hsae} argumentam que modelos supervisionados falham em ataques \textit{zero-day} por dependerem de amostras rotuladas de cada classe de ataque no treinamento. A alternativa proposta é treinar exclusivamente com tráfego benigno e detectar ataques pelo desvio de reconstrução de um \textit{autoencoder}, eliminando a dependência de rótulos de classes maliciosas.

Com base nessas três linhas, são propostos sete cenários adicionais ao \textit{Baseline}, conforme apresentado na \autoref{tab:cenarios_resumo}. O Cenário A reproduzirá o experimento base substituindo o \textit{CICIDS2017} pelo \textit{LycoS-IDS2017}, mantendo os mesmos modelos e variantes de avaliação. Os Cenários B e C aplicarão MDI sobre cada \textit{dataset}, incluindo uma variante com pesos de classe balanceados. Os Cenários D a G investigarão o \textit{Autoencoder} nas quatro combinações entre os dois \textit{datasets} e as duas configurações de \textit{features}. Esse desenho permite isolar o efeito de cada fator e compará-los diretamente.

\begin{table}[H]
\centering
\caption{Resumo dos cenários experimentais planejados}
\label{tab:cenarios_resumo}
\begin{tabular}{llll}
\hline
\textbf{Cenário} & \textbf{\textit{Dataset}} & \textbf{MDI} & \textbf{Modelo(s)} \\
\hline
\textit{Baseline} & \textit{CICIDS2017}    & Não & \textit{Random Forest} + DNN \\
A                 & \textit{LycoS-IDS2017} & Não & \textit{Random Forest} + DNN \\
B                 & \textit{CICIDS2017}    & Sim & DNN \\
C                 & \textit{LycoS-IDS2017} & Sim & DNN \\
D                 & \textit{CICIDS2017}    & Não & \textit{Autoencoder} \\
E                 & \textit{LycoS-IDS2017} & Não & \textit{Autoencoder} \\
F                 & \textit{CICIDS2017}    & Sim & \textit{Autoencoder} \\
G                 & \textit{LycoS-IDS2017} & Sim & \textit{Autoencoder} \\
\hline
\end{tabular}
\end{table}

